\documentclass[../../book.tex]{subfiles}

\begin{document}

\subsection{Unstructured Lattices}\label{subsection:lattices}

We begin our discussion by defining a lattice, its key properties, and assumptions on which lattice-based cryptography relies.
We assume familiarity with standard linear algebra described in \Cref{section:linear-algebra}.

Without further ado, here is the definition of a lattice.

\begin{definition}\label{def:lattice}
    Let $r \leq d$ and let $\mathbf{B} \in \mathbb{R}^{d \times r}$ be some matrix with linearly independent columns $\mathbf{b}_1,\dots,\mathbf{b}_r \in \mathbb{R}^d$.
    We call a \textbf{lattice} an additive group 
    \begin{equation*}
        \mathcal{L}(\mathbf{B}) := \left\{\mathbf{B}\cdot\mathbf{z}: \mathbf{z} \in \mathbb{Z}^r\right\} = \left\{\sum_{i \in [r]}z_i\mathbf{b}_i: z_1,\dots,z_r \in \mathbb{Z}\right\}.    
    \end{equation*}
    We say that the \textbf{rank} of this lattice is $r$.
\end{definition}

Instead of writing $\mathcal{L}(\mathbf{B})$ every time, we will commonly denote lattices by $\Lambda$.

Here is how can geometrically interpret the above definition: if vector space $\mathbb{R}^d$ is typically thought as a set of all $\mathbb{R}$-linear combinations\footnote{We use notation $S+S':=\{s+s': s \in S, s' \in S'\}$ and for $a \in S$, set $aS := \{as: s \in S\}$.} $\sum_{i \in [d]}\mathbf{e}_i\mathbb{R}$ over some basis $\{\mathbf{e}_i\}_{i \in [d]}$, the lattice $\Lambda$ can be thought as all $\mathbb{Z}$-linear combinations $\sum_{i \in [r]}\mathbf{b}_i\mathbb{Z}$ over real basis $\{\mathbf{b}_i\}_{i \in [r]} \subseteq \mathbb{R}^d$.
See \Cref{fig:lattice} for illustration.

\begin{figure}[ht]
  \centering
  \begin{tikzpicture}
    \coordinate (Origin)   at (0,0);
    \coordinate (XAxisMin) at (-2,0);
    \coordinate (XAxisMax) at (5,0);
    \coordinate (YAxisMin) at (0,-2);
    \coordinate (YAxisMax) at (0,5);
    \draw [ultra thick, gray,-stealth] (XAxisMin) -- (XAxisMax);
    \draw [ultra thick, gray,-stealth] (YAxisMin) -- (YAxisMax);

    \clip (-2,-2) rectangle (5cm,5cm); 
    \pgftransformcm{1}{0.6}{0.7}{1}{\pgfpoint{0cm}{0cm}}

    \coordinate (Bone) at (0,2);
    \coordinate (Btwo) at (2,-2);
    \draw[style=help lines, dashed, ultra thick, black!20] (-14,-14) grid[step=2cm] (14,14);

    \foreach \x in {-7,-6,...,7}{
      \foreach \y in {-7,-6,...,7}{
        \node[draw,circle,inner sep=2pt,fill] at (2*\x,2*\y) {};
      }
    }

    \draw [ultra thick,-latex,blue!70!black] (Origin)
        -- (Bone) node [right] {$\mathbf{b}_1$};
    \draw [ultra thick,-latex,blue!70!black] (Origin)
        -- (Btwo) node [below right] {$\mathbf{b}_2$};
    \filldraw[fill=red!70!purple, dashed, very thick, fill opacity=0.3, draw=black] 
        (Origin) -- (Bone) -- ($(Bone) + (Bone) - (Btwo)$) -- ($(Bone)-(Btwo)$) -- cycle;
  
    %\draw [ultra thick,-latex,red!70!purple] (Origin)
    %    -- ($0+1*(Bone)-3*(Btwo)$) node [above left] {$\mathbf{b}_1-3\mathbf{b}_2$};
    \draw [ultra thick,-latex,red!70!purple] (Origin)
        -- ($(Bone)-(Btwo)$) node [above left] {$\mathbf{b}_1-\mathbf{b}_2$};
    \filldraw[fill=blue!70!black, dashed, very thick, fill opacity=0.3, draw=black] 
    (Origin) -- (Bone) -- ($(Bone)+(Btwo)$) -- (Btwo) -- cycle;
    \node at ($0.5*(Bone)+0.5*(Btwo)$) {\footnotesize$\mathcal{P}(\textcolor{blue!70!black}{\mathbf{B}})$};
    \node at ($0.5*(Bone) + 0.5*($(Bone)-(Btwo)$)$) {\footnotesize$\;\mathcal{P}(\textcolor{red!70!purple}{\mathbf{B}'})$};
  \end{tikzpicture}
  \caption{Lattice $\Lambda=\mathcal{L}(\textcolor{blue!70!black}{\mathbf{B}})$ of rank
  $2$ with basis $\textcolor{blue!70!black}{\mathbf{B}} =
  \{\textcolor{blue!70!black}{\mathbf{b}_1},\textcolor{blue!70!black}{\mathbf{b}_2}\}$.
  $\textcolor{blue!70!black}{\mathbf{B}}$ is an example of a ``good'' basis where basis
  vectors are almost orthogonal. The basis $\textcolor{red!70!purple}{\mathbf{B}'} =
  \{\textcolor{blue!70!black}{\mathbf{b}_1},\textcolor{red!70!purple}{\mathbf{b}_1-\mathbf{b}_2}\}$
  of the same lattice $\Lambda=\mathcal{L}(\textcolor{red!70!purple}{\mathbf{B}'})$ is ``worse'' (the vectors seem more parallel and
  fundamental domain $\mathcal{P}(\textcolor{red!70!purple}{\mathbf{B}'})$ is skewed).}
  \label{fig:lattice}
\end{figure}

\begin{remark}\label{remark:lattice-in-integers}
Note that the basis of lattice $\Lambda$ is not required to lie in $\mathbb{Z}^d$, compared to coefficients of the linear combination.
However, in practice they usually do: note that we still can represent basis in memory only in the rational form (that is, $\mathbf{b}_i \in \mathbb{Q}^d$).
With the proper rescaling, representation in $\mathbb{Z}^d$ would suffice.
\end{remark}

\begin{remark}
    Further, we consider \textit{full-rank} lattices only: \textit{i.e.}, for $\Lambda \subseteq \mathbb{R}^d$, its rank is $d$. 
    Indeed, if the rank of the lattice is $r<d$, we can always consider the lattice as a subset of $\mathbb{R}^r$.
\end{remark}

\subsubsection*{Motivation}

The definition above seems not overly complicated and consequently one might wonder what makes lattices useful.
We note that long before cryptographers noticed that lattices are a source of a large number of computationally hard problems, lattices had been first proven to be useful in Mathematics, where they have been used going back to at least to the 18th century. 
In fact, lattices naturally arised in many areas such as number theory or sphere packings (with both being one way or another relevant to lattice and code-based cryptography).
One of the reasons why lattice is a common object to encounter in mathematics is the following lemma.

\begin{lemma}
    Any \emph{discrete} additive subgroup of $\mathbb{R}^d$ is a lattice.
\end{lemma}

The above definition should be clear for the reader except for a term \emph{discrete}.
The set $X \subseteq \mathbb{R}^d$ is called \emph{discrete} if $\delta := \inf_{u \neq v}\|u-v\| > 0$ for $u,v \in X$: in other words, distance between any two points is bounded below by a positive constant $\delta \in \mathbb{R}_{>0}$.

So the lemma states that whenever we have a discrete set $X \subseteq \mathbb{R}^d$ (where every point is isolated) and we have an additive group structure, we can find a basis $\mathbf{b}_1,\dots,\mathbf{b}_d$ such that every element of $X$ can be described as a $\mathbb{Z}$-linear combination of this basis.
This fact is not very hard to prove yet is not immediately obvious from the definitions.

Needless to say, such constructions are frequent in Mathematics. Here is one non-obvious example.

\begin{proposition}
    Let $\mathbb{G}$ be a group generated by $g_1,\dots,g_d \in \mathbb{G}$.
    Define the map $\phi: \mathbb{Z}^d \to \mathbb{G}$ as $(z_1,\dots,z_d) \mapsto z_1g_1 + \dots + z_dg_d$.
    The kernel of this map (that is, all $(z_1,\dots,z_d) \in \mathbb{Z}^d$ such that $\phi(z_1,\dots,z_d) = 0$) is a lattice.
\end{proposition}

\begin{exercise}
    Verify this proposition.
\end{exercise}

As a consequence of observation above, one can prove many algebra theorems (such as \emph{fundamental theorem of finitely generated additive groups}) using theory of lattices.
Specifics are well beyond the scope of this section.
We note that in our further cryptographic applications we will operate primarily over lattices with explicitly generated basis (\Cref{def:lattice}).

\subsubsection{Fundamental Measures}

In this section, we introduce some fundamental characteristics of lattices that will be used throughout the book.

A basis $\mathbf{B}$ of $\Lambda := \mathcal{L}(\mathbf{B})$ is associated with the \textbf{fundamental
domain}\footnote{Sometimes, the interval $[-1/2, 1/2)$ is used
instead of $[0,1)$ to center $\mathcal{P}(\mathbf{B})$ with respect to the origin: in such
case, such set is called a \textit{centered parallelepiped} of $\mathbf{B}$ and we will denote it by $\mathcal{P}_{1/2}(\mathbf{B})$.}:
\begin{equation*}
    \mathcal{P}(\mathbf{B}) \triangleq \left\{ \sum_{i \in [d]} \theta_i\mathbf{b}_i: \theta_i \in [0,1) \right\}.
\end{equation*}

What is special about this parallelepiped is that it represents all elements in the quotient group $\mathbb{R}^d/\Lambda$: any element $\mathbf{x} \in \mathbb{R}^d$ can be uniquely written as $\mathbf{x} = \mathbf{v} + \mathbf{p}$ where $\mathbf{v} \in \Lambda$ and $\mathbf{p} \in \mathcal{P}(\mathbf{B})$.
Moreover, the shape of $\mathcal{P}(\mathbf{B})$ determines the ``quality'' of the basis $\mathbf{B}$:
if $\mathbf{B}$ consists of long, nearly-parallel vectors, the corresponding centered parallelepiped
$\mathcal{P}_{1/2}(\mathbf{B})$ is skewed and has points far from the origin, and $\mathbf{B}$ is typically
considered ``bad,'' as some computational problems become much harder to solve.
In turn, a basis $\mathbf{B}'$ whose centered parallelepiped $\mathcal{P}_{1/2}(\mathbf{B}')$ is non-skewed
and stays close to the origin is considered ``good,'' as $\{\mathbf{b}_i'\}_{i \in [d]}$ much more
closely resembles an orthogonal basis than $\mathbf{B}$ does (and of course, the orthogonal basis is in
a sense ``perfect,'' because it is very ``nice'' to work with).
Again, see \Cref{fig:lattice} for illustration.

Another natural quantity to consider is the volume of such parallelepiped, called \textit{covolume}.

\begin{definition}[Covolume]
    The \textbf{covolume} $\mathsf{Vol}(\Lambda)$ of lattice $\Lambda$ is defined as\footnote{Since $\mathbb{R}^d/\Lambda$ is a smooth manifold, we use $\mathsf{Vol}(\cdot)$ above to denote the natural Lebesgue measure.}:
    \begin{equation*}
        \mathsf{Vol}(\Lambda) \triangleq \mathsf{Vol}(\mathbb{R}^d/\Lambda) = \sqrt{\det \mathbf{B}^{\top}\mathbf{B}}.
    \end{equation*}
\end{definition}

Note that the definition of $\mathcal{P}(\mathbf{B})$ depends on the choice of $\mathbf{B}$ since the basis of $\Lambda$ is generally not unique.
Thus, it might seem that covolume should also depend on the choice of $\mathbf{B}$.
However, one can show the following.

\begin{proposition}
    The \textbf{covolume} $\mathsf{Vol}(\Lambda)$ is an invariant of $\Lambda$:
    that is, it does not depend on the choice of $\mathbf{B}$.
\end{proposition}

Indeed, if $\mathbf{B}_0$ and $\mathbf{B}_1$ are both bases of $\Lambda$, then $\mathbf{B}_1=\mathbf{B}_0\mathbf{U}$ for some unimodular matrix $\mathbf{U} \in \text{GL}_d(\mathbb{Z})$.
In such case one easily observes that $\det \mathbf{B}^{\top}_i\mathbf{B}_i$ is the same for $i \in \{0,1\}$.

Turns out that one of the hardest computational problems is computing the shortest vectors of $\Lambda$.
However, we need some preliminary work to formalize this statement.
Denote $B_{\rho}(\mathbf{x}_0) := \{\mathbf{x} \in \mathbb{R}^d: \|\mathbf{x}-\mathbf{x}_0\|<\rho\}$ --- an open ball in $\mathbb{R}^d$ of radius $\rho$ centered at $\mathbf{x}_0$.
We omit specifying $\mathbf{x}_0$ when it is the origin $\mathbf{0}$.

\begin{definition}[Successive minima]
    The \textbf{$i^{\text{th}}$ minimum $\lambda_i(\Lambda)$} is the radius 
    of the smallest sphere centered at the origin containing $i$ linearly 
    independent lattice vectors:
    \begin{equation*}
        \lambda_i(\Lambda) \triangleq \inf \{\rho \in \mathbb{R}_{\geq 0}: \dim(\text{span}(\Lambda \cap B_{\rho})) \geq i\}.
    \end{equation*}
\end{definition}

In particular, $\lambda_1(\Lambda) = \min_{\mathbf{x} \neq \mathbf{y} \in \Lambda}\|\mathbf{x}-\mathbf{y}\| = \min_{\mathbf{x} \in \Lambda\setminus\{\mathbf{0}\}}\|\mathbf{x}\|$ is called the \textit{first minimum}, while the vector of length $\lambda_1(\Lambda)$ is the shortest vector.
The natural question to ask though is whether it exists at all: after all, in $\mathbb{R}^d$, for instance, it is clearly undefined.
We give an affirmative answer by means of the following proposition.

\begin{proposition}\label{prop:lattice-minima-existence}
    Let $\Lambda \subseteq \mathbb{R}^d$ be a full-rank lattice with successive minima $\lambda_1, \dots, \lambda_d$.
    Then, the following two facts are true:
    \begin{enumerate}
        \item $\lambda_d > \ldots > \lambda_1 > 0$;
        \item there exists a set of linearly independent vectors $\mathbf{v}_1, \dots, \mathbf{v}_d \in \Lambda$ such that
        $\|\mathbf{v}_i\| = \lambda_i$ for all $i \in [d]$.
    \end{enumerate}
    
\end{proposition}

\begin{proof}
    See \cite[Theorem 1.2]{MicGol02}. 
\end{proof}

Note that \Cref{prop:lattice-minima-existence} provides a more general result than what we need. 
Specifically, it implies the shortest vector $\mathbf{v}_1$ of $\Lambda$ exists and has a positive length $\lambda_1(\Lambda)$.

\subsubsection{Lattice Computational Problems}\label{section:lattice-problems}

In this section, we define some problems that are currently believed to be unsolvable in polynomial time (under reasonable choice of parameters). 

\paragraph{Shortest Vector Problem (SVP)}
As previously announced, we start with the \textit{shortest vector problem} (SVP), which requires finding the vector of length $\lambda_1(\Lambda)$.
\begin{definition}[Lattice Generation Algorithm]
    A \textbf{lattice generation algorithm} $\mathsf{LGen}(1^{n})$ is a probabilistic polynomial time (PPT) algorithm that takes the security parameter $1^{n}$ and outputs the lattice basis\footnote{As previously noted in Remark~\ref{remark:lattice-in-integers}, we might assume that $\mathbf{B}$ is given in the integer representation.} $\mathbf{B} \in \mathbb{Z}^{n \times n}$ where $\Lambda=\mathcal{L}(\mathbf{B})$ is of rank $n$.
\end{definition}

\begin{definition}[Shortest Vector Problem]\label{def:svp}
    For a given adversary $\A$ and lattice generation algorithm $\mathsf{LGen}(1^{n})$, define $\underline{\GAME\,\mathsf{SVP}}^{\A}_{\mathsf{LGen}}$ as follows:
    \begin{algoen}
        \item Sample the lattice basis $\mathbf{B} \gets \mathsf{LGen}(1^{n})$ and send to $\A$;
        \item $\A$ outputs some $\mathbf{x} \in \mathbb{Z}^{n} \setminus \{\mathbf{0}\}$ based on $\mathbf{B} \in \mathbb{Z}^{n \times n}$;
        \item \textbf{return} 1 iff $\|\mathbf{B}\mathbf{x}\|=\lambda_1(\Lambda)$.
    \end{algoen}
    We define the $\mathcal{A}$'s \textbf{advantage in solving the shortest
    vector problem} as:
    \begin{equation*}
        \Adv^{\mathsf{SVP}}_{\A,\mathsf{LGen}}(n) := \Pr\left[\underline{\GAME\,\mathsf{SVP}}^{\A}_{\mathsf{LGen}}(n) = 1\right]
    \end{equation*} 
\end{definition}

\begin{definition}[SVP Assumption]
    We say that \textbf{Shortest Vector Problem (SVP) assumption} holds if for 
    any PPT adversary $\mathcal{A}$,
    \begin{equation*}
        \Adv^{\mathsf{SVP}}_{\A,\mathsf{LGen}}(n) = \negl(n).
    \end{equation*}
\end{definition}

\paragraph{Approximate Shortest Vector Problem ($\text{SVP}_{\gamma}$)}

Turns out that the setting of \Cref{def:svp} is too strict: finding the shortest non-zero vector in lattice $\Lambda$ is too hard.
Moreover, one expects that building cryptographic protocols solely on this assumption would be very hard.
So instead, we will require not finding the shortest vector, but ``short-enough'', where the natural measure of ``short'' is ``the small multiple of $\lambda_1(\Lambda)$''.

\begin{definition}[Approximate Shortest Vector Problem]\label{def:appr-svp}
    For a given adversary $\A$ and lattice generation algorithm $\mathsf{LGen}(1^{n})$, define $\underline{\GAME\,\mathsf{SVP}}^{\A}_{\textcolor{red!70!purple}{\gamma},\mathsf{LGen}}$ as follows:
    \begin{algoen}
        \item Sample the lattice basis $\mathbf{B} \gets \mathsf{LGen}(1^{n})$ and send to $\A$;
        \item $\A$ outputs some $\mathbf{x} \in \mathbb{Z}^{n} \setminus \{\mathbf{0}\}$ based on $(\mathbf{B},\textcolor{red!70!purple}{\gamma})$;
        \item \textbf{return} 1 iff $\|\mathbf{B}\mathbf{x}\| \leq \textcolor{red!70!purple}{\gamma}\lambda_1(\Lambda)$.
    \end{algoen}
    We define the $\mathcal{A}$'s \textbf{advantage in solving the approximate
    shortest vector problem} as:
    \begin{equation*}
        \Adv^{\mathsf{SVP}_{\textcolor{red!70!purple}{\gamma}}}_{\A,\mathsf{LGen}}(n) := \Pr\left[\underline{\GAME\,\mathsf{SVP}}^{\A}_{\textcolor{red!70!purple}{\gamma},\mathsf{LGen}}(n) = 1\right]
    \end{equation*} 
\end{definition}


\begin{definition}[Approximate SVP Assumption]\label{def:approximate-svp-assumption}
    We say that \textbf{Approximate Shortest Vector Problem (SVP) assumption}
    with given $\gamma$ holds if for any PPT adversary $\mathcal{A}$,
    \begin{equation*}
        \Adv^{\mathsf{SVP}_{\gamma}}_{\A,\mathsf{LGen}}(n) = \negl(n).
    \end{equation*}
\end{definition}

Typically, $\gamma$ is analyzed as a function of the rank of the lattice $n$.
According to \cite[Section 2.2.5]{KoenWessel25}, for example, vectors of length below $\gamma\lambda_1(\Lambda)$ for $\gamma = \sqrt{n}$ are considered ``very short'' while $\gamma = \mathsf{poly}(n)$ is the most interesting case.
As one expects, the Approximate SVP assumption does not hold for $\gamma=2^{n}$ by means of Lenstra-Lenstra-Lovasz (LLL) algorithm~\cite{LenLenLov82} (even more specifically, LLL solves $\mathsf{SVP}_{\gamma}$ for $\gamma > 2^{(n\log\log n)/\log n}$).

Additionally, one commonly introduces \textbf{(approximate) shortest independent vectors problem ($\text{SIVP}_{\gamma}$)} with almost identical formulation to \Cref{def:appr-svp}, but instead one must find $n$ linearly independent vectors of lattice $\Lambda$ of length at most $\gamma\lambda_{n}(\Lambda)$.
Its hardness is similar to that of $\mathsf{SVP}_{\gamma}$: for instance, algorithm solving $\mathsf{SVP}_{\gamma}$ automatically solves $\text{SIVP}_{\gamma\sqrt{n}}$.

\paragraph{Approximate Closest Vector Problem ($\text{CVP}_{\gamma}$)}\label{section:approximate-cvp}
A somewhat similar problem to SVP is, given some target vector $\mathbf{t} \in \mathbb{Z}^{n}$, find the closest vector to $\mathbf{t}$ on the lattice.
This problem also turns out to be computationally hard.
Similarly to SVP, the raw CVP problem is too strict, so we give formal definition for approximate CVP straight away.

\begin{definition}[Approximate Closest Vector Problem]\label{def:cvp}
    For a given adversary $\A$ and lattice generation algorithm $\mathsf{LGen}(1^{n})$, define $\underline{\GAME\,\mathsf{CVP}}^{\A}_{\gamma,\mathsf{LGen}}$ as follows:
    \begin{algoen}
        \item Sample the lattice basis $\mathbf{B} \gets \mathsf{LGen}(1^{n})$;
        \item Sample target vector $\mathbf{t} \leftarrowS \mathbb{Z}^{n}$ and send $(\mathbf{B},\mathbf{t})$ to $\A$;
        \item $\A$ outputs some $\mathbf{x} \in \mathbb{Z}^{n}$ based on $(\mathbf{B},\mathbf{t},\gamma)$;
        \item \textbf{return} $1$ iff $\|\mathbf{B}\mathbf{x}-\mathbf{t}\| \leq \gamma\|\mathbf{B}\mathbf{y}-\mathbf{t}\|$ for all $\mathbf{y} \neq \mathbf{x}$ and $0$ otherwise.
    \end{algoen}
    We define the $\mathcal{A}$'s \textbf{advantage in solving the approximate closest
    vector problem} as:
    \begin{equation*}
        \Adv^{\text{CVP}_{\gamma}}_{\A,\mathsf{LGen}}(n) := \Pr\left[\underline{\GAME\,\mathsf{CVP}}^{\A}_{\gamma,\mathsf{LGen}}(n) = 1\right]
    \end{equation*} 
\end{definition}

Similarly to~\Cref{def:approximate-svp-assumption}, we say that approximate CVP assumption with the approximate factor $\gamma$ holds if $\Adv^{\text{CVP}_{\gamma}}_{\A,\mathsf{LGen}}(n)=\negl(n)$ for all PPT adversaries $\A$.

\paragraph{Bounded Distance Decoding Problem ($\text{BDD}_{\gamma}$)}\label{section:bdd}

Sometimes it is more convenient to use even a looser notion than $\text{CVP}_{\gamma}$.
Specifically, we now allow the target vector $\mathbf{t}$ to be ``close enough'' to the given lattice $\Lambda$.
Formally, we formulate this problem below (and mark changes in \textcolor{blue!70!purple}{blue}).

\begin{definition}[Bounded Distance Decoding
    ($\text{BDD}_{\gamma}$)]\label{def:bdd} For a given adversary $\A$
    and lattice generation algorithm $\mathsf{LGen}(1^{n})$, define
    $\underline{\GAME\,\mathsf{BDD}}^{\A}_{\gamma,\mathsf{LGen}}$
    as follows:
    \begin{algoen}
        \item Sample the lattice basis $\mathbf{B} \gets \mathsf{LGen}(1^{n})$;
        \item Sample target vector $\mathbf{t} \leftarrowS \mathbb{Z}^{n}$
        \textcolor{blue!70!purple}{with $\mathsf{dist}(\mathbf{t}, \Lambda) <
        \lambda_1(\Lambda)/2\gamma$} and send $(\mathbf{B},\mathbf{t})$ to $\A$;
        \item $\A$ outputs some $\mathbf{x} \in \mathbb{Z}^{n}$ based on
        $(\mathbf{B},\mathbf{t},\gamma)$;
        \item \textbf{return} $1$ iff \textcolor{blue!70!purple}{$\mathbf{x} \in
        \Lambda$ \textbf{and}
        $\|\mathbf{x}-\mathbf{t}\|<\lambda_1(\Lambda)/2\gamma$}, and $0$
        otherwise.
    \end{algoen}
    We define the $\A$'s \textbf{advantage in solving the bounded
    distance decoding problem} as:
    \begin{equation*}
        \Adv^{\text{BDD}_{\gamma}}_{\A,\mathsf{LGen}}(n) := \Pr\left[\underline{\GAME\,\mathsf{BDD}}^{\A}_{\gamma,\mathsf{LGen}}(n) = 1\right].
    \end{equation*} 

    We say that $\text{BDD}_{\gamma}$ holds if
    $\Adv^{\text{BDD}_{\gamma}}_{\A,\mathsf{LGen}}(n) = \negl(n)$ for all PPT $\A$.
\end{definition}

\paragraph{Decisional Shortest Vector Problem ($\mathsf{GapSVP}_{\gamma}$)}\label{section:gap-cvp}

Finally, since we have a \textit{search} version of $\mathsf{SVP}_{\gamma}$, it is reasonable to ask for the \textit{decisional} version of the same problem.
Moreover, turns out that the most current protocols are based on the decisional version of $\mathsf{SVP}_{\gamma}$ (called $\mathsf{GapSVP}_{\gamma}$) rather than on the search one.

\begin{definition}[Decisional Shortest Vector Problem
    ($\mathsf{GapSVP}_{\gamma}$)]\label{def:gap-cvp} For a given adversary
    $\A$ and lattice generation algorithm $\mathsf{LGen}(1^{n})$,
    define
    $\underline{\GAME\,\mathsf{GapSVP}}^{\A}_{\gamma,\mathsf{LGen}}$
    as follows:
    \begin{algoen}
        \item Sample a random bit $b \leftarrowS \{0,1\}$;
        \item If $b=0$, sample $\Lambda$ such that $\lambda_1(\Lambda) \leq 1$,
        otherwise such that $\lambda_1(\Lambda) > \gamma$;
        \item $\A$ outputs some $\widehat{b}$ based on
        $(\Lambda,\gamma)$;
        \item \textbf{return} 1 iff $b=\widehat{b}$.
    \end{algoen}
    We define the $\A$'s \textbf{advantage in solving the decisional
    shortest vector problem} as:
    \begin{equation*}
        \Adv^{\mathsf{GapSVP}_{\gamma}}_{\A,\mathsf{LGen}}(n) := \left|\Pr\left[\underline{\GAME\,\mathsf{GapSVP}}^{\A}_{\gamma,\mathsf{LGen}}(n) = 1\right] - \frac{1}{2}\right|
    \end{equation*} 
    We say that $\mathsf{GapSVP}_{\gamma}$ holds if
    $\Adv^{\mathsf{GapSVP}_{\gamma}}_{\A,\mathsf{LGen}}(n)=\negl(n)$
    for all PPT $\A$.
\end{definition}

\begin{remark}
    A core characteristic of these lattice problems is their presumed resistance not only against classical adversaries but quantum as well. 
    Specifically, currently one of the best attacks against $\mathsf{SVP}_{\gamma}$ is given by BDGL sieve algorithms, requiring time $2^{0.292n+\overline{o}(n)}$ for classical computers and $2^{0.257n+\overline{o}(n)}$ for quantum adversaries (with the lower bound of $2^{0.2075n+\overline{o}(n)}$)~\cite{sieve-complexity}.
    For more comprehensive review, see \cite{KoenWessel25}.
    That said, taking $n \gtrapprox 1024$ typically suffices to ensure $\geq 128$ bits of security.
\end{remark}

\subsection{Short Integer Solution (SIS)}\label{section:sis}

The first problem that serves as a foundation of many subsequent protocols (such as collision-resistant hash functions, commitments, and digital signatures) is the \textit{short integer solution} (SIS) problem.
Informally, the SIS problem, as the name suggests, asks for a \emph{short} non-zero solution $\mathbf{x} \in \mathbb{Z}^m$ to the homogeneous linear equation $\mathbf{A}\mathbf{x} = \mathbf{0}$ over $\mathbb{Z}_q$, where $\mathbf{A} \in \mathbb{Z}_q^{n \times m}$ is a random matrix.
At first glance, this problem has nothing to do with the lattices from \Cref{section:lattice-problems}; we establish the connection later in this section.
 
Without any additional constraints, the equation is easy to solve. Indeed, for instance $\mathbf{x} = (q,0,\dots,0)$ is a trivial non-zero solution. By using Gaussian elimination, one may additionally find many ``less trivially looking'' ones.
That said, we additionally require $\|\mathbf{x}\|_{\infty} \leq \beta$ for some fixed bound $\beta \ll q$, and this requirement is what makes the problem hard.

\paragraph{Notation}
Throughout the rest of this section, $q$ is an (odd) modulus and we represent every $x \in \mathbb{Z}_q$ by its \emph{centered} representative in the range from $-q/2$ to $q/2$.
Consequently, $\|x\|_{\infty}$ for $x \in \mathbb{Z}_q$ denotes the absolute value of this representative (for example, $\|q-1\|_{\infty} = 1$), and naturally let $\|\mathbf{x}\|_{\infty} := \max_i \|x_i\|_{\infty}$ and $\|\mathbf{x}\|_2 := \sqrt{\sum_{i \in [m]} \|x_i\|_{\infty}^2}$ for $\mathbf{x} \in \mathbb{Z}_q^m$.
We call the elements of
\begin{equation*}
    S_{\eta} := \{x \in \mathbb{Z}_q: \|x\|_{\infty} \leq \eta\}
\end{equation*}
\emph{short}, and we implicitly assume $\eta \ll q$ (think of $\eta \leq 4$ and $q \approx 2^{23}$).

We denote the SIS problem with parameters $(n,m,q,\beta)$ by $\mathsf{SIS}_{n,m,q,\beta}$ and give its formal definition below.

\begin{definition}[SIS Assumption]\label{def:sis}
    For a given adversary $\A$ and parameters $(n,m,q,\beta)$, define $\underline{\GAME\,\mathsf{SIS}}^{\A}_{n,m,q,\beta}$ as follows:
    \begin{algoen}
        \item Sample a matrix $\mathbf{A} \leftarrowS \mathbb{Z}_q^{n \times m}$ and send it to $\A$;
        \item $\A$ outputs some $\mathbf{x} \in \mathbb{Z}^m$;
        \item \textbf{return} $1$ iff $\mathbf{x} \neq \mathbf{0}$, $\mathbf{A}\mathbf{x} = \mathbf{0} \pmod{q}$, and $\|\mathbf{x}\|_{\infty} \leq \beta$.
    \end{algoen}
    We define the $\A$'s \textbf{advantage in solving the SIS problem} as
    \begin{equation*}
        \Adv^{\mathsf{SIS}}_{\A}(n) := \Pr\left[\underline{\GAME\,\mathsf{SIS}}^{\A}_{n,m,q,\beta}(n) = 1\right].
    \end{equation*}
    We say that the \textbf{$\mathsf{SIS}_{n,m,q,\beta}$ assumption} holds if $\Adv^{\mathsf{SIS}}_{\A}(n) = \negl(n)$ for any PPT adversary $\A$.
\end{definition}

\begin{remark}
    Some works measure shortness in the $\ell^2$ norm instead of the $\ell^{\infty}$ norm.
    Since $\|\mathbf{x}\|_{\infty} \leq \|\mathbf{x}\|_2 \leq \sqrt{m}\|\mathbf{x}\|_{\infty}$, the two variants differ only by a factor of at most $\sqrt{m}$ in the bound $\beta$.
\end{remark}

\paragraph{Choosing parameters}
The parameters $(n,m,q,\beta)$ are not independent, and the following observations explain which choices make sense~\cite{EPRINT:Peikert15}:
\begin{enumerate}
    \item[\textit{(i)}] For $\beta \geq q$ the problem is obviously trivial: for example, $(q,0,\dots,0)$ is a valid non-zero solution.
    \item[\textit{(ii)}] Increasing $m$ can only make SIS easier: a solution $\mathbf{x}$ for $\mathbf{A}$ extends to the solution $(\mathbf{x},\mathbf{0})$ for $[\mathbf{A}\,\|\,\mathbf{A}']$ with the same norm.
    Similarly, increasing $n$ (the number of equations) can only make SIS harder.
    \item[\textit{(iii)}] A short solution exists as soon as $(\beta+1)^m > q^n$.
    Indeed, the map $\mathbf{x} \mapsto \mathbf{A}\mathbf{x}$ sends $(\beta+1)^m$ vectors $\mathbf{x} \in \{0,\dots,\beta\}^m$ into a set of size $q^n$, so by the pigeonhole principle two distinct vectors $\mathbf{x} \neq \mathbf{x}'$ collide.
    Their difference $\mathbf{x}-\mathbf{x}'$ is a non-zero solution with $\|\mathbf{x}-\mathbf{x}'\|_{\infty} \leq \beta$.
    Typically, one therefore takes $m > n\log q/\log(\beta+1)$, which is much larger than $n$.
\end{enumerate}

\paragraph{SIS as a lattice problem}
To see why SIS is a lattice problem, associate with $\mathbf{A} \in \mathbb{Z}_q^{n \times m}$ the set of all integer solutions:
\begin{equation*}
    \Lambda^{\perp}_q(\mathbf{A}) := \{\mathbf{x} \in \mathbb{Z}^m: \mathbf{A}\mathbf{x} = \mathbf{0} \pmod{q}\}.
\end{equation*}
This set is closed under addition and discrete, so it is a lattice.
Moreover, $q\mathbb{Z}^m \subseteq \Lambda^{\perp}_q(\mathbf{A}) \subseteq \mathbb{Z}^m$, so $\Lambda^{\perp}_q(\mathbf{A})$ has full rank $m$; lattices ``sandwiched'' between $q\mathbb{Z}^m$ and $\mathbb{Z}^m$ are called \textbf{$q$-ary lattices}.
One can also show that $\mathsf{Vol}(\Lambda^{\perp}_q(\mathbf{A})) = q^n$ whenever the rows of $\mathbf{A}$ are linearly independent over $\mathbb{Z}_q$ (which happens with overwhelming probability for random $\mathbf{A}$).

A solution to SIS is by definition a short non-zero vector of $\Lambda^{\perp}_q(\mathbf{A})$.
Therefore, SIS is exactly the approximate shortest vector problem (\Cref{def:appr-svp}) on the $q$-ary lattice $\Lambda^{\perp}_q(\mathbf{A})$. 
With a bit of additional care, one can show that the approximation factor of the corresponding SVP is roughly $\beta/\lambda_1(\Lambda^{\perp}_q(\mathbf{A}))$.

In other words, if we knew how to solve approximate SVP on the lattice $\Lambda^{\perp}_q(\mathbf{A})$, we would be able to break the SIS assumption: SIS is no harder than $\mathsf{SVP}_{\gamma}$ on this particular lattice.
This observation is useful for attacks, but it does not tell us that SIS is hard.
For security we need a reduction in the opposite direction, and, remarkably, it exists: as first shown by Ajtai, solving this \emph{average-case} problem is at least as hard as solving certain lattice problems in the \emph{worst case}.
We state the modern version of this result informally.

\begin{theorem}[Informal, see {\cite[Section 4.1]{EPRINT:Peikert15}}]\label{thm:sis-worst-case}
    For $m = \mathsf{poly}(n)$, any $\beta > 0$, and any sufficiently large $q \geq \beta \cdot \mathsf{poly}(n)$, solving $\mathsf{SIS}_{n,m,q,\beta}$ with non-negligible probability is at least as hard as solving $\mathsf{GapSVP}_{\gamma}$ and $\mathsf{SIVP}_{\gamma}$ on \emph{arbitrary} $n$-dimensional lattices for some $\gamma = \beta \cdot \mathsf{poly}(n)$.
\end{theorem}

In other words, breaking a random SIS instance breaks \emph{every} lattice of dimension $n$, so the protocol designer does not need to worry about sampling ``weak'' instances.

\paragraph{Application: hash functions and commitments}
The most direct application of SIS is a collision-resistant hash function.
For a random key $\mathbf{A} \leftarrowS \mathbb{Z}_q^{n \times m}$ with $2^m > q^n$, define
\begin{equation*}
    \mathsf{H}_{\mathbf{A}}: \{0,1\}^m \to \mathbb{Z}_q^n, \quad \mathbf{x} \mapsto \mathbf{A}\mathbf{x}.
\end{equation*}
The condition $2^m > q^n$ makes $\mathsf{H}_{\mathbf{A}}$ compressing: it maps $m$ bits into $n\log q < m$ bits.
As observed earlier, this also ensures that the equation $\mathbf{Ax}=0$ has at least one short solution, which will be important for proving the following proposition.

\begin{proposition}
    If $\mathsf{SIS}_{n,m,q,1}$ holds, then $\mathsf{H}_{\mathbf{A}}$ is collision-resistant.
\end{proposition}

\begin{proof}
    Suppose an adversary finds $\mathbf{x} \neq \mathbf{x}' \in \{0,1\}^m$ with $\mathbf{A}\mathbf{x} = \mathbf{A}\mathbf{x}'$.
    Then $\mathbf{A}(\mathbf{x}-\mathbf{x}') = \mathbf{0}$, the vector $\mathbf{x}-\mathbf{x}'$ is non-zero, and each of its coordinates lies in $S_1$.
    Hence, $\mathbf{x}-\mathbf{x}'$ is a solution to $\mathsf{SIS}_{n,m,q,1}$ for the challenge matrix $\mathbf{A}$.
\end{proof}

Notice that, contrary to SHA-256, the function $\mathsf{H}_{\mathbf{A}}$ is linear: $\mathsf{H}_{\mathbf{A}}(\mathbf{x}) + \mathsf{H}_{\mathbf{A}}(\mathbf{x}') = \mathsf{H}_{\mathbf{A}}(\mathbf{x}+\mathbf{x}')$.
This is exactly the property that made the Pedersen commitment useful, and it gives the lattice analogue of Pedersen commitments, commonly called the \textbf{Ajtai commitment}.
We define it as follows.

\begin{definition}[Ajtai Commitment]
    Define the \textbf{Ajtai commitment scheme} $\Pi_{\text{Ajtai}}$ over a message space $S_{\delta}^n$ and randomness space $S_{\eta}^m$ as follows:
    \begin{itemize}
        \item $\mathsf{Setup}(1^{\lambda}) \to \mathsf{pp}$: generate two matrices $\mathbf{A} \leftarrowS \mathbb{Z}_q^{\kappa \times n}$ and $\mathbf{B} \leftarrowS \mathbb{Z}_q^{\kappa \times m}$.
        \item $\mathsf{Com}(\mathbf{x} \in S_{\delta}^n; \mathbf{r} \in S_{\eta}^m) \to (\mathbf{c} \in \mathbb{Z}_q^{\kappa})$: to generate a commitment to a short message $\mathbf{x} \in S_{\delta}^n$, the prover selects a blinding factor $\mathbf{r} \leftarrowS S_{\delta}^m$ and outputs $\mathbf{c} \gets \mathbf{Ax} + \mathbf{Br}$.
    \end{itemize}
\end{definition}

Binding follows from SIS exactly as in the proof above: two openings $(\mathbf{x},\mathbf{r}) \neq (\mathbf{x}',\mathbf{r}')$ give the short solution $(\mathbf{x}-\mathbf{x}',\mathbf{r}-\mathbf{r}')$.

\begin{exercise}
    Which SIS instance does this collision solve?
    Specify the matrix and the parameters $(n', m', q', \beta')$ of $\mathsf{SIS}_{n',m',q',\beta'}$ (\Cref{def:sis}).

    \commentgray{Answer: $\mathsf{SIS}_{\kappa,n+m,q,2\max\{\delta,\eta\}}$.}
\end{exercise}

Unfortunately, the hiding property cannot be shown using SIS (namely, we will require the LWE assumption which will be introduced in the next section).
The commitment is additively homomorphic, but only as long as the committed values stay short: adding many commitments increases the norm of the underlying message, and binding is guaranteed only up to the bound $\beta$.
Keeping track of such norm growth is the recurring theme of all lattice-based protocols.

\begin{exercise}
    Let $\mathbf{x}_1,\dots,\mathbf{x}_N \in S_1^n$ be committed with Ajtai commitments $\mathbf{c}_1,\dots,\mathbf{c}_N$.
    What SIS bound $\beta$ is needed so that $\mathbf{c}_1+\dots+\mathbf{c}_N$ is still a binding commitment to $\mathbf{x}_1+\dots+\mathbf{x}_N$? Assume $\eta \geq 1$.
    
    \commentgray{Answer: $\beta=2N\eta$.}
\end{exercise}

\paragraph{Inhomogeneous SIS (ISIS)}
Public-key primitives need a one-way function that maps a short secret to a public value.
The natural candidate is $\mathbf{s} \mapsto \mathbf{A}\mathbf{s}$ for short $\mathbf{s}$, and inverting it is the \textbf{inhomogeneous SIS} problem: given $\mathbf{A} \leftarrowS \mathbb{Z}_q^{n \times m}$ and $\mathbf{t} \in \mathbb{Z}_q^n$, find $\mathbf{x}$ with $\|\mathbf{x}\|_{\infty} \leq \beta$ and $\mathbf{A}\mathbf{x} = \mathbf{t}$.
ISIS is as hard as SIS up to small changes in parameters~\cite{STOC:GenPeiVai08}: intuitively, $\mathbf{t} = \mathbf{A}\mathbf{s}$ for a short random $\mathbf{s}$ is close to uniform when $m$ is large, and a solver that finds a \emph{different} short preimage $\mathbf{x} \neq \mathbf{s}$ yields the SIS solution $\mathbf{x}-\mathbf{s}$.
We will use this one-way function for the Lyubashevsky signature in \Cref{section:lyubashevsky}.

\subsection{Learning With Errors (LWE)}\label{section:lwe}

SIS is well suited for hash functions and commitments, but it is less convenient for encryption: we need a public key that looks random, yet hides a secret that allows decryption.
The \textit{learning with errors} (LWE) problem, introduced by Regev, provides exactly this.

\paragraph{Motivation}
Fix a secret $\mathbf{s} \in \mathbb{Z}_q^n$.
Publishing random linear equations $b_i = \langle \mathbf{a}_i, \mathbf{s} \rangle$ for $\mathbf{a}_i \leftarrowS \mathbb{Z}_q^n$ is useless: after collecting $n$ equations, Gaussian elimination recovers $\mathbf{s}$ in $O(n^3)$ time.
To resolve this issue, LWE perturbs each equation by a small error $e_i$:
\begin{equation*}
    b_i = \langle \mathbf{a}_i, \mathbf{s} \rangle + e_i, \quad e_i \leftarrowS \chi,
\end{equation*}
where $\chi$ is an \emph{error distribution} over short elements of $\mathbb{Z}_q$ (for now, think of the uniform distribution over $S_{\eta}$ for $\eta \leq 4$).
Stacking $m$ such equations yields a matrix $\mathbf{A} \in \mathbb{Z}_q^{m \times n}$ and a vector $\mathbf{b} = \mathbf{A}\mathbf{s} + \mathbf{e}$.
For such a construction, the Gaussian elimination no longer works and no algorithms are known to solve this equation with respect to $\mathbf{s}$.
Let us now formalize what we have so far.

\begin{definition}[LWE Instance]
    For parameters $(n,m,q,\chi)$ and a secret $\mathbf{s} \in \mathbb{Z}_q^n$, the \textbf{LWE instance generation algorithm} $\mathsf{LWEGen}_{n,m,q,\chi}(\mathbf{s})$ samples $\mathbf{A} \leftarrowS \mathbb{Z}_q^{m \times n}$ and $\mathbf{e} \leftarrowS \chi^m$, and outputs $(\mathbf{A}, \mathbf{b} := \mathbf{A}\mathbf{s}+\mathbf{e})$.
\end{definition}

Similarly to Diffie-Hellman assumptions, the LWE problem has two variants: the search version asks to find $\mathbf{s}$, while the decisional version asks to distinguish an LWE instance from a uniformly random pair.
We formalize this below.

\begin{definition}[Search-LWE Assumption]\label{def:slwe}
    For a given adversary $\A$ and parameters $(n,m,q,\chi)$, define $\underline{\GAME\,\mathsf{SLWE}}^{\A}_{n,m,q,\chi}$ as follows:
    \begin{algoen}
        \item Sample the secret $\mathbf{s} \leftarrowS \mathbb{Z}_q^n$ and $(\mathbf{A},\mathbf{b}) \gets \mathsf{LWEGen}_{n,m,q,\chi}(\mathbf{s})$;
        \item $\A$ outputs some $\widehat{\mathbf{s}}$ based on $(\mathbf{A},\mathbf{b})$;
        \item \textbf{return} $1$ iff $\widehat{\mathbf{s}} = \mathbf{s}$.
    \end{algoen}
    We say that the \textbf{Search-LWE assumption} holds if for any PPT $\A$
    \begin{equation*}
        \Adv^{\mathsf{SLWE}}_{\A}(n) := \Pr[\underline{\GAME\,\mathsf{SLWE}}^{\A}_{n,m,q,\chi} = 1] = \negl(n).
    \end{equation*}
\end{definition}

\begin{definition}[Decisional-LWE Assumption]\label{def:dlwe}
    For a given adversary $\A$ and parameters $(n,m,q,\chi)$, define $\underline{\GAME\,\mathsf{DLWE}}^{\A}_{n,m,q,\chi}$ as follows:
    \begin{algoen}
        \item Sample a random bit $b \leftarrowS \{0,1\}$;
        \item If $b=0$, sample $\mathbf{s} \leftarrowS \mathbb{Z}_q^n$ and $(\mathbf{A},\mathbf{u}) \gets \mathsf{LWEGen}_{n,m,q,\chi}(\mathbf{s})$; otherwise, sample $(\mathbf{A},\mathbf{u}) \leftarrowS \mathbb{Z}_q^{m \times n} \times \mathbb{Z}_q^m$;
        \item $\A$ outputs some $\widehat{b}$ based on $(\mathbf{A},\mathbf{u})$;
        \item \textbf{return} $1$ iff $\widehat{b} = b$.
    \end{algoen}
    We say that the \textbf{Decisional-LWE assumption} holds if for any PPT $\A$
    \begin{equation*}
        \Adv^{\mathsf{DLWE}}_{\A}(n) := \left|\Pr\left[\underline{\GAME\,\mathsf{DLWE}}^{\A}_{n,m,q,\chi}(n) = 1\right] - \frac{1}{2}\right| = \negl(n).
    \end{equation*}
\end{definition}

Clearly, if one can find $\mathbf{s}$, one can distinguish (compute $\mathbf{b}-\mathbf{A}\mathbf{s}$ and check whether it is short).
Regev showed that the converse also holds for $q = \mathsf{poly}(n)$, so the two assumptions are equivalent.
The idea is to recover $\mathbf{s}$ one coordinate at a time.
To test whether a guess $g \in \mathbb{Z}_q$ equals $s_1$, sample $\boldsymbol{\Delta} \leftarrowS \mathbb{Z}_q^m$, add it to the first column of $\mathbf{A}$, and add $g\boldsymbol{\Delta}$ to $\mathbf{b}$.
If $g = s_1$, the modified pair is again an LWE instance with the same secret; otherwise, $\mathbf{b}$ receives the extra term $(g-s_1)\boldsymbol{\Delta}$, which is uniform and makes the pair uniform.
The distinguisher tells the two cases apart, and trying all $q$ guesses for each of the $n$ coordinates takes $qn$ calls.

\begin{exercise}
    Check that the modified pair in the argument above is uniformly distributed when $g \neq s_1$ and $q$ is prime.
\end{exercise}

\paragraph{Short secrets}
Most practical schemes sample the secret from the error distribution, $\mathbf{s} \leftarrowS \chi^n$, instead of uniformly over $\mathbb{Z}_q^n$.
This variant is called \emph{short-secret LWE}.
Although it drastically reduces the number of possible secrets, it does not reduce security: short-secret LWE with $m$ samples is at least as hard as standard LWE with $m+n$ samples.
In the rest of the book we use short secrets implicitly.

\paragraph{LWE as a lattice problem}
Similarly to SIS, the matrix $\mathbf{A} \in \mathbb{Z}_q^{m \times n}$ defines the $q$-ary lattice
\begin{equation*}
    \Lambda_q(\mathbf{A}) := \{\mathbf{y} \in \mathbb{Z}^m: \mathbf{y} = \mathbf{A}\mathbf{x} \pmod{q} \; \text{for some} \; \mathbf{x} \in \mathbb{Z}^n\} = \mathbf{A}\mathbb{Z}^n + q\mathbb{Z}^m.
\end{equation*}
The vector $\mathbf{b} = \mathbf{A}\mathbf{s}+\mathbf{e}$ does not lie on $\Lambda_q(\mathbf{A})$, but it is very close to the lattice point $\mathbf{A}\mathbf{s} \bmod q$: the distance is at most $\|\mathbf{e}\|_2 \leq \sqrt{m}\eta$.
Therefore, Search-LWE is an instance of the bounded distance decoding problem (\Cref{def:bdd}): an algorithm solving BDD on $\Lambda_q(\mathbf{A})$ recovers the lattice point $\mathbf{A}\mathbf{s}$, and hence $\mathbf{s}$.
In the other direction, LWE also enjoys a worst-case hardness guarantee, although through a \emph{quantum} reduction.

\begin{theorem}\label{thm:lwe-worst-case}
    For $m = \mathsf{poly}(n)$, a modulus $q \leq 2^{\mathsf{poly}(n)}$, and a (discrete) Gaussian error distribution\footnote{We have not yet defined discrete Gaussian distributions, so think of them as a particularly useful type of short distributions.} $\chi$ of width $\alpha q \geq 2\sqrt{n}$ for $\alpha \in (0,1)$, solving $\mathsf{DLWE}_{n,m,q,\chi}$ is at least as hard as \emph{quantumly} solving $\mathsf{GapSVP}_{\gamma}$ and $\mathsf{SIVP}_{\gamma}$ on arbitrary $n$-dimensional lattices for some $\gamma = \widetilde{O}(n/\alpha)$.
\end{theorem}

\paragraph{Application: public-key encryption}
The decisional form of LWE makes encryption almost immediate.
The public key is an LWE instance $(\mathbf{A},\mathbf{b} = \mathbf{A}\mathbf{s}+\mathbf{e})$, which by DLWE looks uniformly random.
To encrypt a bit $\mu$, the sender takes a random subset-sum of the rows of $(\mathbf{A},\mathbf{b})$ and adds $\mu \cdot \lceil q/2 \rfloor$ to the last coordinate.
The owner of $\mathbf{s}$ cancels the $\mathbf{A}$-part, leaving the message plus a small error.

\begin{algoproto}[ht]
    \small
    \centering
    \begin{mdframed}[linecolor=black, linewidth=1pt,
                     skipabove=10pt, skipbelow=10pt,
                     innerleftmargin=5pt, innerrightmargin=5pt]
        \begin{minipage}[t]{0.3\linewidth}
            \colorbox{oc-blue-1}{$\mathsf{KeyGen}(1^{\lambda}) \to (\mathsf{pk}, \mathsf{sk})$:}
            \begin{algoen}
                \item $\mathbf{A} \leftarrowS \mathbb{Z}_q^{m \times n}$;
                \item $\mathbf{s} \leftarrowS \mathbb{Z}_q^n$, $\mathbf{e} \leftarrowS \chi^m$;
                \item $\mathbf{b} \gets \mathbf{A}\mathbf{s} + \mathbf{e}$;
                \item $\mathsf{pk} \gets (\mathbf{A},\mathbf{b})$;
                \item $\mathsf{sk} \gets \mathbf{s}$;
                \item \textbf{return} $(\mathsf{pk}, \mathsf{sk})$.
            \end{algoen}
        \end{minipage}
        \hfill
        \begin{minipage}[t]{0.32\linewidth}
            \colorbox{oc-blue-1}{$\mathsf{Enc}(\mathsf{pk}, \mu \in \{0,1\}) \to c$:}
            \begin{algoen}
                \item $\mathbf{x} \leftarrowS \{0,1\}^m$;
                \item $\widetilde{\mathbf{a}} \gets \mathbf{A}^{\top}\mathbf{x}$;
                \item $\widetilde{b} \gets \langle \mathbf{b}, \mathbf{x} \rangle + \mu\lceil q/2 \rfloor$;
                \item \textbf{return} $c=(\widetilde{\mathbf{a}}, \widetilde{b})$.
            \end{algoen}
        \end{minipage}
        \hfill
        \begin{minipage}[t]{0.35\linewidth}
            \colorbox{oc-blue-1}{$\mathsf{Dec}(\mathsf{sk}, c) \to \mu$:}
            \begin{algoen}
                \item Parse $(\widetilde{\mathbf{a}}, \widetilde{b}) \gets c$;
                \item $\delta \gets \widetilde{b} - \langle \mathbf{s}, \widetilde{\mathbf{a}} \rangle$;
                \item \textbf{return} $1$ iff $\|\delta\|_{\infty} \geq q/4$.
            \end{algoen}
        \end{minipage}
    \end{mdframed}
    \caption{Key generation, encryption, and decryption in Regev's scheme $\Pi_{\mathsf{PKE}}^{\text{LWE}}$ (with $m \geq (n+1)\log q$).}
    \label{fig:lwe-pke}
\end{algoproto}

\textbf{Correctness.} Expanding $\delta$ gives
\begin{equation*}
    \delta = \mathbf{x}^{\top}(\mathbf{A}\mathbf{s}+\mathbf{e}) + \mu\lceil q/2 \rfloor - \mathbf{s}^{\top}\mathbf{A}^{\top}\mathbf{x} = \textcolor{oc-green-7}{\langle \mathbf{e},\mathbf{x} \rangle} + \mu\lceil q/2 \rfloor.
\end{equation*}
Since $\mathbf{x}$ is binary, the error term \textcolor{oc-green-7}{$\langle \mathbf{e},\mathbf{x} \rangle$} is a sum of at most $m$ values from $S_{\eta}$, so $\|\langle \mathbf{e},\mathbf{x} \rangle\|_{\infty} \leq m\eta$.
If $m\eta < q/4$, then $\delta$ is close to $0$ for $\mu = 0$ and close to $q/2$ for $\mu = 1$, and decryption is always correct.

\textbf{Security.} The scheme is IND-CPA secure under DLWE, and the argument has two steps.
First, by DLWE we may replace $\mathbf{b}$ with a uniformly random vector without the adversary noticing.
Second, once $(\mathbf{A},\mathbf{b})$ is uniform, the leftover hash lemma (this is where $m \geq (n+1)\log q$ is used) implies that $(\mathbf{A}^{\top}\mathbf{x}, \langle \mathbf{b},\mathbf{x} \rangle)$ is statistically close to uniform, so the ciphertext carries no information about $\mu$.

\begin{remark}\label{remark:lwe-inefficiency}
    Notice how inefficient this scheme is: the public key takes $m(n+1)\log q$ bits, and each \emph{bit} of the message costs a ciphertext of $(n+1)\log q$ bits.
    For $n = 512$ and $q \approx 2^{12}$ this is already about $6000$ bits of ciphertext per bit of plaintext.
    Fixing this inefficiency is the main motivation for the structured lattices of the next section.
\end{remark}

\subsection{Structured Lattices*}\label{section:structured-lattices}

The unstructured SIS and LWE problems have a single practical drawback: size.
The matrix $\mathbf{A}$ consists of $nm$ elements of $\mathbb{Z}_q$, which is hundreds of kilobytes for secure parameters, and every multiplication by $\mathbf{A}$ costs $O(nm)$ operations.
Sampling $\mathbf{A}$ from a public seed\footnote{Since $\mathbf{A}$ is uniformly random, all parties can expand it from a short seed with an extendable-output function such as SHAKE. Every practical scheme does this, so the size of $\mathbf{A}$ itself rarely matters.} removes the storage cost of $\mathbf{A}$, but the other components (such as $\mathbf{b}$ in LWE, or the ciphertexts in \Cref{remark:lwe-inefficiency}) remain large, and so does the computation.

All modern lattice-based schemes (including the NIST standards ML-KEM and ML-DSA) solve this problem by replacing the ring $\mathbb{Z}_q$ with a polynomial ring $\mathcal{R}_q$.
A single element of $\mathcal{R}_q$ is described by $n$ numbers, yet multiplication by it acts as an $n \times n$ matrix.
In this section, we define this ring, explain why it is chosen the way it is, and restate SIS and LWE over it.

\subsubsection{Cyclotomic Rings}

Fix $n$ to be a power of two and define the rings
\begin{equation*}
    \mathcal{R} := \mathbb{Z}[T]/(T^n+1), \quad \mathcal{R}_q := \mathbb{Z}_q[T]/(T^n+1).
\end{equation*}
Elements of $\mathcal{R}_q$ are polynomials $f = f_0 + f_1T + \dots + f_{n-1}T^{n-1}$ with coefficients in $\mathbb{Z}_q$.
We add them coefficient-wise, and we multiply them as usual polynomials followed by the reduction $T^n = -1$.
The polynomial $T^n+1$ is the $2n$-th \emph{cyclotomic} polynomial, which is why $\mathcal{R}$ is called a \textbf{cyclotomic ring}.

\paragraph{Multiplication as a matrix}
The \emph{coefficient embedding} identifies $f \in \mathcal{R}$ with its coefficient vector $\mathsf{vec}(f) := (f_0,\dots,f_{n-1})^{\top} \in \mathbb{Z}^n$.
Multiplication by a fixed $a \in \mathcal{R}$ is a linear map on coefficient vectors, so it corresponds to some matrix $\mathsf{rot}(a) \in \mathbb{Z}^{n \times n}$ with
\begin{equation*}
    \mathsf{vec}(a \cdot f) = \mathsf{rot}(a) \cdot \mathsf{vec}(f) \quad \text{for all} \; f \in \mathcal{R}.
\end{equation*}
The $j$-th column of $\mathsf{rot}(a)$ is $\mathsf{vec}(T^j a)$, and multiplying by $T$ shifts the coefficients up by one position while the top coefficient wraps around with a minus sign (since $T^n = -1$).
Hence, $\mathsf{rot}(a)$ is the \emph{anti-circulant} matrix
\begin{equation*}
    \mathsf{rot}(a) = \begin{bmatrix}
        a_0 & -a_{n-1} & \dots & -a_1 \\
        a_1 & a_0 & \dots & -a_2 \\
        \vdots & \vdots & \ddots & \vdots \\
        a_{n-1} & a_{n-2} & \dots & a_0
    \end{bmatrix}.
\end{equation*}

\begin{exercise}
    For $n=2$, compute $(1+2T)(3-T)$ in $\mathbb{Z}[T]/(T^2+1)$ and check that the result equals $\mathsf{rot}(1+2T)\cdot(3,-1)^{\top}$.
\end{exercise}

Therefore, a single ring element describes a structured $n \times n$ matrix, and a matrix over $\mathcal{R}_q$ of size $k \times \ell$ describes a structured matrix over $\mathbb{Z}_q$ of size $kn \times \ell n$ built from blocks $\mathsf{rot}(a_{ij})$.
This is the source of all efficiency gains: storing the structured matrix takes $k\ell n$ instead of $k\ell n^2$ elements of $\mathbb{Z}_q$.

\subsubsection{Ideal and Module Lattices*}

\begin{remark}
    The material presented below is fairly complicated. 
    The key idea of employing structured lattices is the following: instead of using vectors from $\mathbb{Z}_q^n$, we use polynomials from the $2n$-th cyclotomic ring $\mathcal{R}_q$.
    The section below explains why such a substitution makes sense.
\end{remark}

Let us see where lattices arise when considering $\mathcal{R}$.
Recall that an \emph{ideal} $\mathcal{I} \subseteq \mathcal{R}$ is an additive subgroup closed under multiplication by elements of $\mathcal{R}$; the simplest example is the principal ideal $(g) := \{gf: f \in \mathcal{R}\}$.
The image $\mathsf{vec}(\mathcal{I}) \subseteq \mathbb{Z}^n$ of an ideal is a lattice, called an \textbf{ideal lattice}.
For a principal ideal, $\mathsf{vec}((g)) = \mathcal{L}(\mathsf{rot}(g))$, so the $n$ coefficients of $g$ describe the entire basis.

\begin{example}
    Let $\mathcal{R} = \mathbb{Z}[T]/(T^2+1)$ and $\mathcal{I} = (1+T)$.
    Every element of $\mathcal{I}$ has the form $(1+T)(\alpha+\beta T) = (\alpha-\beta) + (\alpha+\beta)T$, so
    \begin{equation*}
        \mathsf{vec}(\mathcal{I}) = \mathcal{L}\left(\begin{bmatrix} 1 & -1 \\ 1 & 1 \end{bmatrix}\right).
    \end{equation*}
\end{example}

Ideal lattices are closed under the rotation $\mathsf{vec}(f) \mapsto \mathsf{vec}(Tf)$, and this extra symmetry is both their strength and their weakness.
For the approximation factors used in cryptography, the best known algorithms for SVP on ideal lattices still take time $2^{\Omega(n)}$, even on quantum computers.
However, the structure does make some problems easier: for instance, $\mathsf{GapSVP}_{\gamma}$ is easy on ideal lattices for small polynomial factors, and quantum algorithms find mildly short vectors (with $\gamma = 2^{\widetilde{O}(\sqrt{n})}$) in polynomial time~\cite{JACM:CraLeoWes21}.
None of these attacks affects practical schemes, but they motivate a more conservative choice.

\paragraph{Module lattices}
A \emph{module} over $\mathcal{R}$ is the analogue of a vector space in which the scalars come from the ring $\mathcal{R}$ instead of a field.
All we need are submodules of $\mathcal{R}^d$ generated by $d$ vectors, and their images under the coefficient embedding applied component-wise,
\begin{equation*}
    \mathsf{vec}(f_1,\dots,f_d) := (\mathsf{vec}(f_1),\dots,\mathsf{vec}(f_d)) \in \mathbb{Z}^{nd},
\end{equation*}
which are called \textbf{module lattices} of rank $d$.
For a matrix $\mathbf{A} \in \mathcal{R}^{k \times \ell}$ we similarly define $\mathsf{rot}(\mathbf{A}) \in \mathbb{Z}^{kn \times \ell n}$ as the block matrix with blocks $\mathsf{rot}(A_{ij})$, which preserves the key property $\mathsf{vec}(\mathbf{A}\mathbf{f}) = \mathsf{rot}(\mathbf{A})\mathsf{vec}(\mathbf{f})$.
Norms extend in the natural way: $\|\mathbf{f}\|_{\infty} := \max_i \|f_i\|_{\infty}$ and $\|\mathbf{f}\|_2 := \|\mathsf{vec}(\mathbf{f})\|_2$.

Module lattices interpolate between the two previous worlds:
\begin{itemize}
    \item for $n = 1$, we have $\mathcal{R} = \mathbb{Z}$, and module lattices are ordinary unstructured lattices in $\mathbb{R}^d$;
    \item for $d = 1$, module lattices are ideal lattices in $\mathbb{R}^n$;
    \item for $n > 1$ and $d > 1$, we get lattices of dimension $nd$ that are structured only in $n \times n$ blocks.
\end{itemize}
No known attack exploits the module structure for rank $d \geq 2$ better than generic lattice algorithms, and the quantum attacks on ideal lattices do not extend to higher ranks.

\subsubsection{Module-SIS and Module-LWE}

We now restate SIS and LWE over $\mathcal{R}_q$.
It is convenient to state the module versions directly, since the ring versions are special cases.
The short distribution $\chi$ now samples each coefficient of an element of $\mathcal{R}_q$ independently from a short distribution over $\mathbb{Z}_q$; 
we abuse notation to denote $S_{\eta} := \{f \in \mathcal{R}_q: \|f\|_{\infty} \leq \eta\}$.

\begin{definition}[Module-SIS Assumption]\label{def:msis}
    For a given adversary $\A$ and parameters $(\mathcal{R},q,k,\ell,\beta)$, define $\underline{\GAME\,\mathsf{MSIS}}^{\A}_{q,k,\ell,\beta}$ as follows:
    \begin{algoen}
        \item Sample a matrix $\mathbf{A} \leftarrowS \mathcal{R}_q^{k \times \ell}$ and send it to $\A$;
        \item $\A$ outputs some $\mathbf{x} \in \mathcal{R}^{\ell}$;
        \item \textbf{return} $1$ iff $\mathbf{x} \neq \mathbf{0}$, $\mathbf{A}\mathbf{x} = \mathbf{0}$ over $\mathcal{R}_q$, and $\|\mathbf{x}\|_{\infty} \leq \beta$.
    \end{algoen}
    We say that the \textbf{$\mathsf{MSIS}_{q,k,\ell,\beta}$ assumption} holds if the probability that this game outputs $1$ is negligible for any PPT adversary $\A$.
\end{definition}

\begin{definition}[Module-LWE Assumption]\label{def:mlwe}
    For parameters $(\mathcal{R},q,k,\ell,\chi)$, a \textbf{Module-LWE instance} is a pair $(\mathbf{A}, \mathbf{b} := \mathbf{A}\mathbf{s}+\mathbf{e})$ for $\mathbf{A} \leftarrowS \mathcal{R}_q^{k \times \ell}$, $\mathbf{s} \leftarrowS \chi^{\ell}$, and $\mathbf{e} \leftarrowS \chi^{k}$.
    Then,
    \begin{enumerate}
        \item The \textbf{Search Module-LWE assumption} ($\mathsf{SMLWE}_{q,k,\ell,\chi}$) states that no PPT adversary recovers $\mathbf{s}$ from $(\mathbf{A},\mathbf{b})$ with non-negligible probability.
        \item The \textbf{Decisional Module-LWE assumption} ($\mathsf{DMLWE}_{q,k,\ell,\chi}$) states that no PPT adversary distinguishes $(\mathbf{A},\mathbf{b})$ from a uniform pair in $\mathcal{R}_q^{k \times \ell} \times \mathcal{R}_q^{k}$ with non-negligible advantage, exactly as in \Cref{def:dlwe}.
    \end{enumerate}
\end{definition}

The \textbf{Ring-SIS} and \textbf{Ring-LWE} problems are the special cases with $k = 1$ (a single equation $\langle \mathbf{a},\mathbf{x} \rangle = 0$ in $\mathcal{R}_q$) and $k = \ell = 1$ (a single sample $b = as+e$), respectively.
Via the map $\mathsf{rot}$, a Module-SIS instance is simply an ordinary SIS instance $\mathsf{rot}(\mathbf{A}) \in \mathbb{Z}_q^{kn \times \ell n}$ with a structured matrix, and similarly for Module-LWE.
Both problems enjoy worst-case to average-case reductions from hard problems on module lattices~\cite{EPRINT:LyuPeiReg12, DCC:AdeDam15}, analogous to \Cref{thm:sis-worst-case} and \Cref{thm:lwe-worst-case}.

\paragraph{Ring-LWE-based Encryption}
Ring-LWE resolves the inefficiency from \Cref{remark:lwe-inefficiency}.
The idea is to encode $n$ message bits $\mu_0,\dots,\mu_{n-1}$ as the coefficients of $\mu := \sum_i \mu_iT^i \in \mathcal{R}_q$ and encrypt all of them at once.
\begin{algoproto}[ht]
    \small
    \centering
    \begin{mdframed}[linecolor=black, linewidth=1pt,
                     skipabove=10pt, skipbelow=10pt,
                     innerleftmargin=5pt, innerrightmargin=5pt]
        \begin{minipage}[t]{0.25\linewidth}
            \colorbox{oc-blue-1}{$\mathsf{KeyGen}(1^{\lambda})$:}
            \begin{algoen}
                \item $a \leftarrowS \mathcal{R}_q$;
                \item $s, e \leftarrowS \chi$;
                \item $b \gets as + e$;
                \item $\mathsf{pk} \gets (a, b)$;
                \item $\mathsf{sk} \gets s$;
                \item \textbf{return} $(\mathsf{pk}, \mathsf{sk})$.
            \end{algoen}
        \end{minipage}
        \hfill
        \begin{minipage}[t]{0.35\linewidth}
            \colorbox{oc-blue-1}{$\mathsf{Enc}(\mathsf{pk}, \boldsymbol{\mu} \in \{0,1\}^n) \to c$:}
            \begin{algoen}
                \item $\mu \gets \sum_{i \in (n)} \mu_i T^i$;
                \item $r, e_1, e_2 \leftarrowS \chi$;
                \item $u \gets ar + e_1$;
                \item $v \gets br + e_2 + \lceil q/2 \rfloor \mu$;
                \item \textbf{return} $c = (u, v)$.
            \end{algoen}
        \end{minipage}
        \hfill
        \begin{minipage}[t]{0.35\linewidth}
            \colorbox{oc-blue-1}{$\mathsf{Dec}(\mathsf{sk}, c) \to \boldsymbol{\mu}$:}
            \begin{algoen}
                \item Parse $(u, v) \gets c$;
                \item $\delta \gets v - su$;
                \item $\mu_i \gets 1$ iff $\|\delta_i\|_{\infty} < q/4$ for each $i \in (n)$;
                \item \textbf{return} $(\mu_0,\dots,\mu_{n-1})$.
            \end{algoen}
        \end{minipage}
    \end{mdframed}
    \caption{Key generation, encryption, and decryption in the Ring-LWE-based scheme $\Pi_{\mathsf{PKE}}^{\text{RLWE}}$ over $\mathcal{R}_q = \mathbb{Z}_q[T]/(T^n+1)$, where $\chi$ is a short distribution over $S_{\eta} \subseteq \mathcal{R}_q$ and $\delta_i$ denotes the $i$-th coefficient of $\delta$.}
    \label{fig:rlwe-pke}
\end{algoproto}

\textbf{Correctness.} Recall that $b = as + e$.
Since $\mathcal{R}_q$ is commutative, the terms $asr$ and $sar$ coincide, and we get
\begin{align*}
    \delta &= v - su \\
    &= (as + e)r + e_2 + \lceil q/2 \rfloor \mu - s(ar + e_1) \\
    &= \textcolor{oc-green-7}{er + e_2 - se_1} + \lceil q/2 \rfloor \mu.
\end{align*}
By the bound $\|fg\|_{\infty} \leq n\|f\|_{\infty}\|g\|_{\infty}$, every coefficient of the error \textcolor{oc-green-7}{$er + e_2 - se_1$} has absolute value at most $2n\eta^2 + \eta$.
Hence, decryption recovers every bit $\mu_i$ whenever $2n\eta^2 + \eta < q/4$.

\textbf{Security.} Security follows from Decisional Ring-LWE, applied twice: once to replace $b$ with a uniform element, and once to replace $(u,v)$, which is itself a pair of Ring-LWE samples with secret $r$.

The ciphertext now consists of two ring elements, \textit{i.e.}, $2n\log q$ bits for $n$ message bits, instead of $n(n+1)\log q$ bits in Regev's scheme.
Encryption costs a few multiplications in $\mathcal{R}_q$, each taking $O(n\log n)$ time with the NTT.
\end{document}